\documentclass[a4j, 12pt]{jarticle}
\usepackage{graphicx}
\usepackage{titling}
\usepackage{amsmath}
\usepackage{float}
\usepackage{listings}
\usepackage{booktabs}
\usepackage{siunitx}
\usepackage{makeidx}
\usepackage{amssymb}
\usepackage{url}
\usepackage{subcaption}
\renewcommand{\lstlistingname}{コード}
\graphicspath{{./assets/eps}}
\author{クリスティアン　ハルジュノ\thanks{釧路工業高等専門学校情報工学科5年情報21番}}
\date{\today}
\begin{document}
\begin{titlepage}
  \centering
  \vspace*{3cm}

  {\fontsize{25pt}{36pt}\selectfont\bfseries グラフ・ネットワークプログラム} \\[1cm]  % Main title
  {\Large 釧路工業高等専門学校情報工学科5年・情報工学実験II} \\[1cm]
  {\large 氏名：クリスティアン・ハルジュノ} \\[1cm]
  {\large 出席番号：21} \\[1cm]
  {\large 学籍番号：234071} \\[5cm]
  {\Large 提出日：\today} \\
  \vfill
\end{titlepage}
\tableofcontents
\newpage
\setcounter{page}{1}
\section{初めに}
グラフネットワーク理論は, 対象間の関係をモデル化するために数学的構造を解析する研究である. 専門的で限られた分野のように見えるが, グラフネットワーク理論は現代の計算機科学において大きな役割を果たしている. 特に, 地域間のインターネット通信を最適化する上で重要である. インターネットの基盤となるサーバーファームは世界中に点在している. 最適化が行われなければ, 自宅からこれらのサーバーへアクセスするには物理的距離の大きさにより多大な時間を要することになる. グラフ理論を用いることで, 情報が通過する最適経路を見出し, インターネットアクセスの高速化を実現することが可能となる.\\

本実験では, ネットワーク理論の基礎を検討し, グラフネットワーク内でノードを探索する複数の手法を比較する. また, 大規模なノードグラフに存在するクリークの数を推定するために, グラフ計算アルゴリズムを最適化する手法についても検討する.
\section{背景}\label{sec:background}
\subsection{グラフ探索アルゴリズム}\label{sec:graph-traversal}
グラフデータ構造を探索するアルゴリズムはいくつか存在する. 本実験では, 代表的なグラフ探索アルゴリズムの中でも特に広く用いられている深さ優先探索 (Depth-First Search, DFS) と幅優先探索 (Breadth-First Search, BFS) の二つを用いてグラフを探索する.

\subsubsection{課題１: 深さ優先探索 (Depth-First Search, DFS)}\label{sec:dfs}
DFSは木あるいはグラフのデータ構造を走査または探索するためのアルゴリズムである. このアルゴリズムは根から探索を開始し, ある枝を可能な限り深く進んだ後に親ノードへ戻り, 他の枝を探索するという手順で動作する. DFSの概念自体は古くから知られ研究されてきたが, Tarjan (1972) がこのアルゴリズムを形式化し, さまざまな問題への応用を示した\cite{tarjan1972dfs}.\\

DFSは幅よりも深さを優先する探索手法である. 最も基本的な形態 (古典的反復DFS) では, 探索対象のノードを保持するためにスタック (後入れ先出し, LIFO) のデータ構造を用いる. アルゴリズムはまず根ノードから探索を開始する. このノードをスタックに積み込み, 取り出した後に訪問済みとして記録する. その後, このノードの子ノードをすべてスタックに積み込み, スタックが空になるまで同じ処理を繰り返す\cite{tarjan1972dfs}. 深さ優先探索アルゴリズムの可視化は図\ref{figure:dfs-example}に示される. \\

\begin{figure}[H]
  \centering
  \includegraphics[width=\textwidth]{dfs-example.eps}
  \caption{深さ優先探索の手順 (DFS) の例}\label{figure:dfs-example}
\end{figure}

しかし, StackOverflowのコミュニティ投稿によれば, ノードをスタックから取り出した後に訪問済みとする方式では, グラフがサイクルを含む場合に同一ノードが複数回スタックに積まれる危険がある. この問題を回避するため, 本研究では修正版である反復DFS (Modified Iterative DFS) を用いる. この手法ではノードをスタックに積み込む前に訪問済みとして記録する. これにより, ノードが再度スタックに積み込まれることを防ぐことができる\cite{kalinichenko2013dfs}.

\subsubsection{課題２: 幅優先探索 (Breadth-First Search, BFS)}\label{sec:bfs}
深さ優先探索 (DFS) と同様に, 幅優先探索 (Breadth-First Search, BFS) アルゴリズムもグラフ・ネットワーク構造の走査あるいは探索に用いられる. DFSが一つの枝を可能な限り深く探索してから次の枝に移るのに対し, BFSはまずノードの子をすべて探索し, その後で次の階層へ進む. DFSがスタック (後入れ先出し, LIFO) を用いるのに対し, BFSはキュー (先入れ先出し, FIFO) を用いて探索対象のノードを管理する\cite{moore1959shortest}.\\

BFSもDFSと同様に根ノードから探索を開始する. 根ノードのすべての子ノード (根から直接導かれるノード) を訪問済みとして記録し, それらをキューに入れる. 次にキューの先頭からノードを取り出し, その子ノードを探索して訪問済みと記録し, キューへ追加する. この処理をキューが空になり探索すべきノードがなくなるまで繰り返す\cite{moore1959shortest}. 幅優先探索アルゴリズムの可視化は図\ref{figure:bfs-example}に示される.\\

\begin{figure}[H]
  \centering
  \includegraphics[width=\textwidth]{bfs-example.eps}
  \caption{幅優先探索の手順 (BFS) の例}\label{figure:bfs-example}
\end{figure}

元のDFSアルゴリズムを修正してサイクルを含むグラフでノードを再訪問しないようにしたのとは対照的に, BFSアルゴリズムはそのままの形でノードの再訪問を回避できる. これは, ノードをキューに追加する前に訪問済みとして記録するためである. この仕組みにより, 各ノードは一度しかキューに入らず, サイクルを含むグラフにおいても無限ループを防ぐことができる\cite{moore1959shortest}.

\subsection{全域木 (Spanning Tree)}\label{sec:spanning-tree}
グラフ探索中に訪問したノードを記録すると, 得られる構造は全域木と呼ばれる. 全域木とは, 元のグラフ内の全ての頂点を含む部分グラフであり, 接続されておりサイクルを含まない木であり, 頂点数よりちょうど1つ少ない辺を持つ\cite{west2001introduction}. 隣接グラフから全域木への変換の例は, 図\ref{figure:spanning-tree-transformation}に示されている.\\

\begin{figure}[H]
  \centering
  \includegraphics[width=\textwidth]{spanning-tree-transformation.eps}
  \caption{隣接グラフから全域木への変換の例}\label{figure:spanning-tree-transformation}  
\end{figure}

グラフは探索の方法によって複数の全域木を持つことがある. DFSやBFSのような異なる探索方法は, 異なる特徴を持つ全域木を生成する. DFS はできるだけ深く探索してからバックトラックするため, 得られる全域木は非常に高さがあり細長くなることがある. 一方, BFS は次のレベルに移動する前に全ての子ノードを探索するため, 得られる全域木は低く広い形になる\cite{west2001introduction}.



\subsection{課題３: 連結成分数}\label{sec:connected-components}
一般に, グラフは大きく2種類に分類される: 連結グラフと非連結グラフである. 連結グラフとは, 任意の2頂点間に経路が存在するグラフを指す. この場合, 任意の頂点から出発して全ての頂点に到達することが可能である\cite{diestel2017graph}. 一方で, 非連結グラフとは, 少なくとも2頂点間に経路が存在しないグラフを指す. この場合, 始点から探索しても到達できない頂点が存在する\cite{west2001introduction}. このような状況により, グラフ内部には「連結成分」が形成される. 非連結グラフは複数の連結成分を持つ場合がある. 連結成分の個数は「連結成分数」と呼ばれる. 例えば, 図\ref{figure:connected-2}では2つの連結成分が存在し, 連結成分数は2となる.\\

\begin{figure}[H]
\centering
\includegraphics[width=0.8\textwidth]{connected-2.eps}
\caption{2つの連結成分を持つ非連結グラフ}\label{figure:connected-2}
\end{figure}

グラフにおける連結成分数を求める方法には, Union-Find, 隣接行列, スペクトルグラフ理論など複数の手法が存在するが, 最も容易で実装しやすい方法はDFSやBFSといったグラフ探索アルゴリズムを用いる方法である. 連結成分数を計算する際には, グラフ内の全ての頂点を順に確認し, 未訪問であれば探索を行う. 探索に用いるアルゴリズムはDFSでもBFSでも問題なく, いずれも最終的には全ての頂点を訪問する. 前章\ref{sec:graph-traversal}で述べたように, 実装されるアルゴリズムは訪問後に頂点を「訪問済み」として記録する. 探索を実行するたびに連結成分数を1増加させる. 探索アルゴリズムは単一の連結成分内の全ての頂点を訪問するため, 最初の探索終了後にさらに探索を実行する場合, そのグラフは非連結であり, 新たな連結成分として数えられる. この処理を全ての頂点が訪問済みになるまで繰り返す. 最終的な連結成分数は探索を実行した回数に等しい\cite{cormen2009introduction}.

\subsection{課題４: 最大クリーク}\label{sec:clique}
\subsubsection{クリークの定義}\label{sec:clique-introduction}
クリークとは, 無向グラフ内の頂点の部分集合で, 集合内の全ての頂点が互いに直接接続されているものを指す. 言い換えると, クリークとはより大きなグラフの一部である完全グラフである\cite{diestel2017graph}. サイズ $k$ のクリーク（$k$-クリークとも呼ばれる）は, 部分集合内の全ての頂点の組が辺によって接続されている $k$ 個の頂点の部分集合である. グラフ内で可能な最大のクリークは最大クリークと呼ばれ, そのサイズはグラフ$G$のクリーク数$\omega(G)$と呼ばれる.\\

例えば, 図\ref{figure:clique-example}では, 各グラフが異なるサイズのクリークを表している. 各グラフはそれぞれ 3-クリーク, 4-クリーク, 5-クリークである\cite{kclique_percolation}. 各グラフは完全グラフであり, そのグラフ内の全てのノードが互いに接続されている. 本実験では, 与えられたグラフ内で最大または最も大きいクリークを見つけることが課題である.

\begin{figure}[H]
  \centering
  \includegraphics[width=\textwidth]{clique-size-examples.eps}
  \caption{クリークの例}\label{figure:clique-example}
\end{figure}

\subsubsection{クリーク問題の計算複雑性}\label{sec:clique-complexity}
クリーク問題は計算複雑性理論における NP-完全問題として知られている \cite{garey1979computers}. NP-完全問題は, 提案された解が多項式時間 ($O(n^k)$, 定数 k に対して) で検証可能な難しい問題のクラスである. 問題が以下の条件を満たす場合, NP-完全と定義される\cite{cormen2009introduction}:
\begin{itemize}
  \item 答えが「はい」または「いいえ」である決定問題である.
  \item 「はい」の解は多項式時間（短時間）で検証可能である.
  \item 「はい」の解の正確性は多項式時間で確認可能である.
  \item 問題は他の任意のNP問題をシミュレートまたは表現可能であり, 多項式時間で検証可能である.
\end{itemize}
\section{実験方法}
\subsection{ハードウェアと環境}
ハードウェアとしてはAppleシリコンのプロセッサに基づいて実験を行う. 
\begin{description}
  \item[機種名] Apple Macbook Pro (M3, 2023年モデル)
  \item[プロセッサ] Apple M3チップ（8コア：高性能コア4＋高効率コア4）
  \item[GPU] 10コア統合型GPU
  \item[メモリ] 16GBユニファイドメモリ
  \item[ストレージ] 1TB SSD
  \item[OS] macOS Sequoia 15.5
  \item[アーキテクチャ] ARM64（Appleシリコン） 
\end{description}
また, 作成した実験プログラムはC言語で書き, GCCコンパイラでコンパイルされる. 
\begin{verbatim}
  Apple clang version 17.0.0 (clang-1700.0.13.5)
  Target: arm64-apple-darwin24.5.0
  Thread model: posix
\end{verbatim}
コンパイルを効率化するため, GCCコンパイラをMakeと同時に利用する. 
\begin{verbatim}
  GNU Make 3.81
  Copyright (C) 2006  Free Software Foundation, Inc.
  This is free software; see the source for copying conditions.
  There is NO warranty; not even for MERCHANTABILITY or 
  FITNESS FOR A PARTICULAR PURPOSE.
  This program built for i386-apple-darwin11.3.0
\end{verbatim}
\subsection{実験対象データ}\label{sec:experimental-data}
本研究で行う全ての実験は, 隣接リストで表現された無向グラフを利用する. 各実験では, ノード数が異なる様々なサイズのグラフが与えられる. 実験で使用されるグラフは以下の通りである.

\subsubsection{課題1と課題2: DFSとBFSの実装}
DFSとBFSの実装は, ノード数と密度が異なる6種類の無向グラフに対してテストされる. これらのグラフは隣接行列で表現され, 密度に基づいて2種類に分類される. 密グラフはノード数に対して辺の数が多く, 疎グラフは相対的に辺の数が少ない. 実験で使用されるグラフの詳細は表\ref{tab:adjacency_matrices} にまとめられている.

\begin{table}[H]
  \centering
  \caption{隣接行列の説明}
  \begin{tabular}{|c|c|c|}
    \hline
    \textbf{ファイル名} & \textbf{ノード数} & \textbf{グラフの種類} \\
    \hline
    search\_100d.dat & 100 & 密グラフ \\
    search\_100s.dat & 100 & 疎グラフ \\
    search\_500d.dat & 500 & 密グラフ \\
    search\_500s.dat & 500 & 疎グラフ \\
    search\_1000d.dat & 1000 & 密グラフ \\
    search\_1000s.dat & 1000 & 疎グラフ \\
    \hline
  \end{tabular}\label{tab:adjacency_matrices}
\end{table}
\subsubsection{課題3: 連結成分数}
連結成分数の実験では, 隣接行列で表現された1つの無向グラフを使用する. このグラフは2000を超えるノードを含み, 辺と連結成分の数は未知である. グラフは \texttt{search\_2000.dat} というファイルで与えられている.

\subsubsection{課題4: 最大クリーク}
最大クリーク探索アルゴリズムは, 隣接行列で表現された5種類の無向グラフに対してテストされる. グラフはサイズや密度が異なり, ノード数は30から300までの範囲にわたる. 実験で使用されるグラフの詳細は表 \ref{tab:clique_graphs} にまとめられている.

\begin{table}[H]
  \centering
  \caption{クリーク探索に使用されるグラフの説明}
  \begin{tabular}{|c|c|c|}
    \hline
    \textbf{ファイル名} & \textbf{ノード数} & \textbf{グラフの種類} \\
    \hline
    clique\_30.dat & 30 & 密グラフ \\
    clique\_50.dat & 50 & 密グラフ \\
    clique\_100.dat & 100 & 密グラフ \\
    clique\_200.dat & 200 & 密グラフ \\
    clique\_300.dat & 300 & 密グラフ \\
    \hline
  \end{tabular}\label{tab:clique_graphs}
\end{table}

\subsubsection{データの整形とパース}
本実験で使用される全てのグラフは, 隣接行列として整形されたテキストファイルで与えられる. 例えばグラフが50ノードを含む場合, 隣接行列は$50 \times 50$の行列として表現される. これは探索において小さな課題をもたらす. あるノードが他のノードに接続されているかを確認するためには, 全てのノードを反復処理して接続の有無を確認する必要がある. その結果, 計算量は$O(n^2)$となり, 大規模なグラフに対しては非効率になる可能性がある. 探索処理を最適化するために, 各隣接行列は事前に処理され, 各ノードがポインタを通じて他のノードにリンク可能な構造体として表現される隣接リストに変換される.\\
本実験で使用される各グラフノードの構造体はコード\ref{lst:graph-node-struct} に示されている.

\begin{lstlisting}[caption=グラフノードの構造,label=lst:graph-node-struct]
struct Node{
  int id;
  bool visited;
  int connected_nodes_count;
  Node **connected_nodes;
};
\end{lstlisting}

現在のノードを追跡するために, id が置かれ, ノードが訪問済みかどうかを示すブール値も記録されている. 接続されているノードの数も構造体に保存されており, 接続ノードを反復処理する際に効率的である. 最後に, 他のノードへの接続を表現するために, ノードへのポインタ配列へのポインタが使用されている. これにより, 全てのノードを反復処理する必要なく, 接続されているノードへ直接アクセスすることが可能となり, グラフ探索を効率的に行うことができる.\\

探索の結果として得られる全域木も, ポインタによってリンクされたノードとして表現される. コード\ref{lst:graph-node-struct}と同様に, 全域木の各ノードはidと隣接ノードへのポインタのリストを含む. 唯一の違いは, 親ノードへのポインタが存在する点である. 全域木ノードの完全な構造はコード\ref{lst:spanning-tree-node-struct}に示されている.

\begin{lstlisting}[caption={全域木ノードの構造},label=lst:spanning-tree-node-struct]
struct TreeNode{
  int id;
  int connected_nodes_count;
  TreeNode **connected_nodes;
  TreeNode *parent_node;
};  
\end{lstlisting}
\subsection{スタックとキューの実装}\label{sec:stack-queue}
DFSとBFSのアルゴリズムを実装するためには, スタックとキューを使用する必要がある\cite{cormen2009introduction}. 標準のC言語ライブラリにはスタックやキューの組み込みデータ構造が存在しないため, 本研究ではこれらのデータ構造をリンクリストを用いて一から実装する.\\

C言語における標準配列は固定サイズであり, 動的にサイズを変更することができない. この制約により, 要素の追加や削除に応じて動的に拡張や縮小を行うスタックやキューには不向きである. 一方, リンクリストは動的にメモリを割り当てることができ, 必要に応じてサイズを拡張または縮小することが可能である. 配列は連続したメモリ領域に割り当てられるが, リンクリストのノードはメモリ上で分散していても, ポインタによって順序構造を維持できる. 各リンクリストノードはデータと次のノードへのポインタを含む\cite{cormen2009introduction}\cite{goodrich2016data}.

\subsubsection{スタック}
スタックは本質的にLIFO (Last In First Out)型のデータ構造である. スタックに最後に追加された要素が最初に取り出される.　章\ref{sec:stack-queue}で述べたように, 本実験におけるスタックは単方向リンクリストを用いて実装される. スタックの構造体はコード\ref{lst:stack-struct}に示されている.
\begin{lstlisting}[caption=スタックの構造,label=lst:stack-struct]
struct StackNode{
  Node *graph_node; // グラフノードへのポインタ
  TreeNode *tree_parent; // 全域木ノードへのポインタ
  StackNode *next;
};
struct Stack{
  StackNode *head;
  int len;
};
\end{lstlisting}
スタックは2層構造を用いて実装される. 1つ目の構造体 \texttt{StackNode} はスタック内の各ノードを表す. スタックの長さとスタックの先頭を管理するために, 2つ目の構造体である \texttt{Stack} が使用される. この実装では, リンクリストの先頭（最初の要素）がスタックのトップを表す. 新しい要素はリンクリストの先頭に追加され, 同じ位置から削除される.
\newpage
\subsubsection{キュー}
スタックとは対照的に, キューは FIFO (First In First Out) 型のデータ構造であり, 最初に追加された要素が優先され, 最初に取り出される. スタックの実装と同様に, キューも単方向リンクリストを用いて実装される. キューの構造体はコード \ref{lst:queue-struct} に示されている.

\begin{lstlisting}[caption=キューの構造,label=lst:queue-struct]
struct QueueNode{
  Node *graph_node; // グラフノードへのポインタ
  TreeNode *tree_parent; // 全域木ノードへのポインタ
  QueueNode *next;
};
struct Queue{
  QueueNode *head;
  QueueNode *tail;
  int len;
};
\end{lstlisting}
\subsection{課題１と課題２: DFSとBFSの実装}
先に章\ref{sec:dfs}および章\ref{sec:bfs}で述べたように, 実装した両方の探索アルゴリズムはノードをスタックまたはキューに入れる前に訪問済みとしてマークするように改変している. この修正によって, サイクルを含むグラフにおいてエンキューされたノードやプッシュされたノードが再訪されることを防げる. 両アルゴリズムの実装は再帰ではなく反復を用いており, 大規模なグラフにおけるスタックオーバーフローの可能性を回避している.

\subsubsection{DFSの実装}
DFSは章\ref{sec:stack-queue}で述べたスタックのデータ構造を用いて実装する. DFSは幅よりも深さを優先的に探索するため, アルゴリズムは以下のように動作する:
\begin{enumerate}
  \item 空のスタックを生成する. 根ノードを訪問済みにマークし, それをスタックにプッシュする.
  \begin{lstlisting}[caption=全域木とスタックの初期化,label=lst:init-queue-tree-2]
TreeNode *root_tree = create_tree_node(
                          root_node->id,
                          NULL
                      );
Stack *stack = init();
root_node->visited = true;
push(stack, root_node, root_tree);
  \end{lstlisting}
  \item スタックが空でない限り, 以下を繰り返す:
  \begin{enumerate}
    \item スタックからノードを1つポップする. このノードが現在のノードとなる.
    \begin{lstlisting}[caption=スタックからノード取り出す,label=lst:pop-stack]
StackNode *current_stack_node = pop(stack);
    \end{lstlisting}
    \item 現在のノードの未訪問の子ノードごとに, それを訪問済みにマークし, 対応する木ノードを生成して現在の木ノードに接続し, スタックにプッシュする.
    \begin{lstlisting}[caption=本ノードの各子ノードをスタックに入れる,label=lst:push-stack-nodes]
Node *children = current_node->connected_nodes[i];
if (children->visited == false){
  children->visited = true;
  TreeNode *child_tree_node = create_tree_node(
                                children->id, 
                                current_tree_node
                              );
  add_child(current_tree_node, child_tree_node);
  push(stack, children, child_tree_node);
}
    \end{lstlisting}
  \end{enumerate}
\end{enumerate}
DFS探索によって得られた全域木は返され, さらに解析を行うことができる.
\subsubsection{BFSの実装}
DFSの実装とは対照的に, BFSは章\ref{sec:stack-queue}で述べたキューのデータ構造を用いる. BFSは深さよりも幅を優先するため, アルゴリズムは以下のように動作する:
\begin{enumerate}
  \item 根の全域木を初期化し, 空のキューを生成する. 根ノードを訪問済みにマークし, それをキューにエンキューする.
  \begin{lstlisting}[caption=全域木とキューの初期化,label=lst:init-queue-tree]
TreeNode *root_tree = create_tree_node(
                          root_node->id,
                          NULL
                      );
Queue *q = init_queue();
root_node->visited = true;
enqueue(q, root_node, root_tree);
  \end{lstlisting}
  \item キューが空でない限り, 以下を繰り返す:
  \begin{enumerate}
    \item キューからノードを1つデキューする. このノードが現在のノードとなる.
    \begin{lstlisting}
    QueueNode *current_queue_node = dequeue(q);
    \end{lstlisting}
    \item 現在のノードの未訪問の子ノードごとに, それを訪問済みにマークし, 対応する木ノードを生成して現在の木ノードに接続し, キューにエンキューする.
    \begin{lstlisting}[caption=本ノードの各子ノードをキューに入れる,label=lst:dequeue-queue]
child_graph->visited = true;
TreeNode *child_tree = create_tree_node(
                          child_graph->id,
                          current_tree_node
                      );
add_child(current_tree_node, child_tree);
enqueue(q, child_graph, child_tree);
    \end{lstlisting}
  \end{enumerate}
\end{enumerate}
BFS探索によって得られた全域木は返され, さらに解析を行うことができる.
\subsubsection{全域木の特徴}
DFSおよびBFSの両方から得られた全域木は, 特性を抽出するために解析できる. 本実験では, 得られた全域木を用いて以下を計算する\cite{west2001introduction}:

\begin{description}
  \item[深さ] 根ノードから最も遠い葉ノード（子を持たないノード）までの距離.
  \item[葉ノード数] 子を持たないノードの数.
  \item[最大分岐数] 任意のノードが持つ子ノードの最大数.
  \item[対称性] 得られた全域木がどの程度均衡しているかを示す指標. 値が高いほど木は均衡している.
  \item[直径] 得られた全域木における任意の2ノード間の最長経路.
\end{description}


\subsection{課題３: 連結成分数}
グラフにおける連結成分数の計算方法は, 章\ref{sec:graph-traversal}で述べたDFSまたはBFSアルゴリズムを用いることができる. 本実験で実装した走査アルゴリズムは, 単一の成分に含まれるノードを訪問済みとしてマークする.  
連結成分数を数えるアルゴリズムの実装をコード\ref{lst:connected-components}に示す.

\begin{lstlisting}[caption=連結成分数の計算,label=lst:connected-components]
  reset_nodes(node_list, node_count);
  int connected_nodes_count = 0;
  for (int i = 0; i < node_count; i++){
    if (node_list[i].visited == false){
      TreeNode *tree_node = build_spanning_tree_bfs(
                                node_list
                            );
      free_tree(tree_node);
      connected_nodes_count++;
    }
  }
  return connected_nodes_count;
\end{lstlisting}

\texttt{node\_list}はグラフ内の全ノードをID順に格納した配列である. 配列内の各ノードは, 接続されたノードへのポインタを保持する. この配列を順に走査し, 未訪問のノードを確認する. 各未訪問ノードに対してBFS (またはDFS) を実行する. 走査アルゴリズムはノードを訪問済みにマークするため, 同じ連結成分に含まれる全ノードは訪問済みとなる.\\
走査を1回行うごとに連結成分数を1増加させる. この処理は, グラフ内のすべてのノードが訪問されるまで繰り返される.


\subsection{課題４: 最大クリーク}
\subsubsection{最大クリーク探索には}\label{sec:clique-search-problem}
グラフ内で最大クリークを求める問題はNP完全問題として知られている. これは多項式時間内で解ける既知のアルゴリズムが存在しないことを意味する. 最大クリークを求める最も基本的な方法はブルートフォース手法である. ブルートフォースではノードの全ての組合せを調べ, それらがクリークを形成するかどうかを確認する. 小規模なグラフであれば現実的であるが, グラフのサイズが大きくなるにつれて計算量は急激に増大する.

ノード数が与えられた場合に調べるべき部分集合の数の可視化を図\ref{figure:node-count-search}に示す. 部分集合の総数はノード数$n$に対して$2^n$で表される. この式は$n$要素から成る集合における全ての部分集合（空集合と集合自身を含む）の数を表す. ノード数が増加すると部分集合の数は指数関数的に増加し, 最大クリーク探索に必要な比較回数は急激に増加する\cite{rosen2019discrete}.

\begin{figure}
  \centering
  \includegraphics[width=\linewidth]{node-count-search.eps}
  \caption{ノード数が与えられた場合に, 調べるべき部分集合数の可視化.}\label{figure:node-count-search}
\end{figure}

ブルートフォースは小さなノード数に対しては単純で分かりやすい手法であるが, ノード数が300に達すると現実的ではなくなる. 例えば, ノード数が300の場合, 部分集合の数は$2^{300}$であり, これはおよそ$1.07 \times 10^{90}$となる. この数は天文学的に大きく, 計算に膨大な時間を要する.
\subsubsection{最大クリーク探索の実装}
グラフにおける最大クリークを求めるために, すべてのノード部分集合がクリークを形成するかを確認する単純なブルートフォースあるいはバックトラッキングアルゴリズムを実装する. 各クリークが確認されるたびに, そのサイズを測定し, 発見された最大のものを記録する. 実装手順は以下の通りである\cite{bron1973algorithm}:
\begin{enumerate}
  \item 関数 \texttt{extend\_clique} は, グラフにおける最大クリークを探索するために設計された再帰的バックトラッキングアルゴリズムである. 
  \item パラメータには以下が含まれる:
    \begin{itemize}
      \item \texttt{clique}: 現在クリークを形成している頂点集合,
      \item \texttt{clique\_size}: 現在クリークに含まれている頂点数,
      \item \texttt{candidates}: クリークを拡張できる可能性のある頂点集合,
      \item \texttt{candidate\_size}: 候補集合のサイズ,
      \item \texttt{node\_count}: グラフ全体の頂点数.
    \end{itemize}
  \item 候補が残っていない場合, 関数は現在のクリークサイズを返す.
  \item 候補集合に含まれる各頂点 $v$ に対して:
    \begin{itemize}
      \item $v$ を一時的にクリークに追加する.
      \item 更新されたクリーク内の全ての頂点と接続している頂点のみを含む新しい候補集合を構築し, クリークの性質を保持する.
      \item アルゴリズムは, 減少した候補集合を用いて \texttt{extend\_clique} を再帰的に呼び出し, クリークをさらに拡張することを試みる.
      \item すべての再帰分岐で発見された最大クリークサイズを記録する.
    \end{itemize}
  \item すべての可能性を探索した後, 関数は発見された最大クリークサイズを返す.
\end{enumerate}
\subsubsection{貪欲彩色アルゴリズム}
章\ref{sec:clique-search-problem}で述べたように, ブルートフォース手法は単純であるが, 大規模なグラフにおいては計算コストの面で大きな制約がある. これを緩和するために, 100ノードおよび300ノードのグラフデータセットに対して貪欲彩色アルゴリズムを実装した. 貪欲彩色アルゴリズムは, グラフの頂点に色を割り当てるヒューリスティック手法である. グラフを彩色することで, 最大クリーク探索時に考慮すべき候補頂点の数を大幅に削減できる\cite{kumlander2005practical}.\\

\begin{figure}[H]
  \centering
  \includegraphics[width=0.8\linewidth]{greedy-coloring.eps}
  \caption{貪欲彩色アルゴリズムの可視化}\label{figure:greedy-coloring}
\end{figure}

貪欲彩色アルゴリズムは, 図\ref{figure:greedy-coloring}のように可視化できる. このアルゴリズムの目的は, 隣接する頂点が同じ色にならないように頂点に色を割り当てることである. Lavrov (2024)によれば, グラフにおける最大クリークのサイズは, グラフを彩色するのに必要な最小の色数（クロマティック数）以下である\cite{lavrov2024lecture25}. 利用可能な候補に対して貪欲彩色アルゴリズムを適用することで, 最大クリークのサイズの上界を得ることができる. これにより, 探索時に考慮すべき候補頂点の数を大幅に削減できる.\\
\newpage
\section{実験結果}
\subsection{課題１: DFSの実装}
表\ref{table:search-results-dfs}において, ファイル名の末尾に$s$が付与されているものは疎グラフを, $d$が付与されているものは密グラフを表す. 全てのデータセットにおいて, 全域木の深さは疎グラフと密グラフの間で大きく異なる. 疎グラフは全域木がより深くなる傾向を示しており, これはノード間の接続が希薄であることを示唆している. また, 葉ノードの数が多いことから, 多くのノードが全域木の最外層に位置していることが分かる. 表\ref{table:leaf-percentage-dfs}を見ると, 100ノードの疎グラフを除き, ほとんどのグラフにおいてノードの90\%以上が葉ノードであることが確認できる.
\begin{table}[H]
  \centering
  \caption{DFSによる各データセットの全域木の特性}\label{table:search-results-dfs}
  \resizebox{\textwidth}{!}{%
    \begin{tabular}{lcccccc}
    \hline
    ファイル名 & ノード数 & 深さ & 葉ノード数 & 最大分岐数 & 対称性 & 直径 \\
    \hline
    search\_100s.dat  & 100  & 23 & 73   & 11  & 0.002490037 & 23 \\
    search\_100d.dat  & 100  & 5  & 96   & 66  & 0.009185571 & 5 \\
    search\_500s.dat  & 500  & 37 & 457  & 46  & 0.000380834 & 37 \\
    search\_500d.dat  & 500  & 6  & 495  & 344 & 0.001820973 & 6 \\
    search\_1000s.dat & 1000 & 48 & 947  & 107 & 0.000196260 & 48 \\
    search\_1000d.dat & 1000 & 9  & 992  & 606 & 0.000850535 & 9 \\
    \hline
    \end{tabular}%
  }
\end{table}

最大分岐因子を見ると, その値は疎グラフと密グラフの間で大きく異なる. 密グラフは疎グラフと比較して分岐因子が高くなる傾向を示し, ノード間の接続がより多いことを示している. 一方で, 全てのデータセットにおいて対称性の値は低く, 0.1\%を超えるグラフは存在しなかった. これは生成された全域木が平衡ではないことを示している. また, 分岐因子とは逆に, 疎グラフは密グラフと比較して直径の値が大きくなる傾向を示し, 疎グラフがより接続性の低い構造であることを裏付けている.
\begin{table}[h]
  \centering
  \caption{DFSによって得られた全域木における葉ノードの割合.}\label{table:leaf-percentage-dfs}
  \begin{tabular}{cccc}
  \hline
  ノード数 & グラフタイプ & 葉ノード数 & 葉ノード割合 \\
  \hline
  100   & 疎 (s)   & 73   & 73\%  \\
  100   & 密 (d)   & 96   & 96\%  \\
  500   & 疎 (s)   & 457  & 91\%  \\
  500   & 密 (d)   & 495  & 99\%  \\
  1000  & 疎 (s)   & 947  & 95\%  \\
  1000  & 密 (d)   & 992  & 99\%  \\
  \hline
  \end{tabular}
\end{table}

\subsection{課題２: BFSの実装}

表\ref{table:search-results-bfs}において, ファイル名の末尾が$s$のものは疎グラフを, $d$のものは密グラフを表す. 全てのデータセットにおいて, 全域木の深さは3から5段階と比較的小さい値に収まっており, グラフのノードが良く接続されていることを示している. また, 葉ノード数が多く, 多くのノードが全域木の最外層に位置していることが分かる. 表\ref{table:leaf-percentage-bfs}を見ると, 100ノード疎グラフを除き, ほとんどのグラフにおいてノード数の90\%以上が葉ノードとなっている.

\begin{table}[H]
  \centering
  \caption{BFSによる各データセットの全域木の特性}\label{table:search-results-bfs}
  \begin{tabular}{lccccccr}
  \hline
  ファイル名 & ノード数 & 深さ & 葉ノード数 & 最大分岐数 & 対称性 & 直径 \\
  \hline
  search\_100s.dat  & 100  & 5 & 70  & 14  & 0.009071578 & 7 \\
  search\_100d.dat  & 100  & 3 & 94  & 66  & 0.009716511 & 4 \\
  search\_500s.dat  & 500  & 4 & 453 & 49  & 0.001929467 & 5 \\
  search\_500d.dat  & 500  & 3 & 493 & 344 & 0.001918439 & 4 \\
  search\_1000s.dat & 1000 & 3 & 946 & 107 & 0.000964235 & 4 \\
  search\_1000d.dat & 1000 & 3 & 993 & 606 & 0.000942066 & 4 \\
  \hline
  \end{tabular}
\end{table}
最大分岐数に着目すると, 疎グラフと密グラフの間で大きく異なることが分かる. 密グラフは疎グラフに比べて分岐数が高く, ノード間の接続が多いことを示している. 対称性の値は全てのデータセットで低く, 0.2\%を超えるグラフは存在せず, 得られた全域木があまり均衡していないことを示している. 直径は4から7の範囲に収まっており, グラフのノードが良く接続されているという観察結果をさらに支持している.
\begin{table}[h]
  \centering
  \caption{BFSによって得られた全域木における葉ノードの割合.}\label{table:leaf-percentage-bfs}
  \begin{tabular}{cccc}
  \hline
  ノード数 & グラフタイプ & 葉ノード数 & 葉ノード割合 \\
  \hline
  100   & 疎 (s)   & 70   & 70\%  \\
  100   & 密 (d)   & 94   & 94\%  \\
  500   & 疎 (s)   & 453  & 91\%  \\
  500   & 密 (d)   & 493  & 99\%  \\
  1000  & 疎 (s)   & 946  & 95\%  \\
  1000  & 密 (d)   & 993  & 99\%  \\
  \hline
  \end{tabular}
\end{table}

\subsection{課題３: 連結成分数}
実装の正確性を確認するために, 課題1および課題2で用いた1000ノードの密グラフに対して連結成分アルゴリズムを実行した. その結果を, 章\ref{sec:experimental-data}で前述した2000ノードのグラフの結果と比較した. 結果を表\ref{table:connected-components-results}に示す.
\begin{table}[H]
  \centering
  \caption{各データセットにおけるDFSおよびBFS探索による連結ノード数と割合}\label{table:connected-components-results}
  \begin{tabular}{lccc}
  \hline
  ファイル名 & 総ノード数 & DFS連結ノード数 & BFS連結ノード数 \\
  \hline
  search\_1000d.dat & 1000 & 1    & 1 \\
  search\_2000.dat  & 2000 & 1686 & 1686\\
  \hline
  \end{tabular}
\end{table}
既知の全域連結グラフに対してアルゴリズムを実行した場合, 期待通り連結成分数は1となった. 一方, より大規模なグラフ (search\_2000.dat) に対して実行した場合, 連結成分数は1686となり, このグラフが完全には連結しておらず, 複数の非連結成分を含むことが確認された. DFSおよびBFSの両実装は同一の結果を出力し, 探索手法の違いによる影響がないことから, 実装の正確性が確認された.
\subsection{課題４: 最大クリーク}
\subsubsection{アルゴリズム確認}
バックトラッキングアルゴリズムの実装が正しいか検証するために, 最大クリークサイズが既知の5ノードからなる小規模テストグラフを作成した. 隣接行列は式\ref{eq:test-adjacency-matrix}に示す.

\begin{equation}\label{eq:test-adjacency-matrix}
  \begin{bmatrix}
  0 & 1 & 1 & 0 & 0 \\
  1 & 0 & 1 & 1 & 0 \\
  1 & 1 & 0 & 1 & 1 \\
  0 & 1 & 1 & 0 & 1 \\
  0 & 0 & 1 & 1 & 0
  \end{bmatrix}
\end{equation}

行列\ref{eq:test-adjacency-matrix}に基づき, 各ノードがどのように接続されているかを表\ref{table:matrix-connections}で確認できる. グラフ中の最大クリークは\{1, 2, 3\}, \{2, 3, 4\}, \{3, 4, 5\}であり, それぞれ完全に接続された部分グラフを形成している. 各クリークに含まれるノード数の最大値は3であるため, 最大クリークサイズは3と判定できる.
\begin{table}[H]
\centering
\caption{各ノードの接続関係}\label{table:matrix-connections}
\begin{tabular}{c|c}
\hline
\textbf{ノード} & \textbf{接続ノード} \\
\hline
1 & 2, 3 \\
2 & 1, 3, 4 \\
3 & 1, 2, 4, 5 \\
4 & 2, 3, 5 \\
5 & 3, 4 \\
\hline
\end{tabular}
\end{table}


実装したバックトラッキングアルゴリズムを実行した結果は, コード\ref{backtracking-test-results}に示す.
\begin{lstlisting}[caption={模擬グラフの実験結果.}, label=backtracking-test-results]
> readnode test_3.dat
File opened successfully.
Biggest Clique: 3
\end{lstlisting}
\subsubsection{ブルートフォース手法のみ}\ref{tab:backtracking-only-results}
すべてのデータセットに対してバックトラッキングアルゴリズムを実行した結果, 最大クリークサイズはそれぞれ\texttt{clique\_30}で8, \texttt{clique\_50}で9, \texttt{clique\_70}で11となった. 30ノードグラフおよび50ノードグラフにおける最大クリークサイズ探索は, いずれも1秒未満で終了した. しかし, 70ノードグラフにおける探索時間は大幅に増加し, アルゴリズムの実行は約22秒で終了した. 
\begin{table}[H]
\centering
\caption{異なるグラフサイズにおける最大クリークサイズと計算時間（バックトラッキングアルゴリズム）}\label{tab:backtracking-only-results}
\begin{tabular}{l|c|r}
\hline
\textbf{グラフサイズ} & \textbf{最大クリークサイズ} & \textbf{計算時間 (s)} \\
\hline
30  & 8  & 0.084 \\
50  & 9  & 0.286 \\
70  & 11 & 21.586 \\
100 & $-$  & $-$ \\
300 & $-$  & $-$ \\
\hline
\end{tabular}
\end{table}

しかし、100 ノードおよび 300 ノードのグラフに対してバックトラッキングアルゴリズムを実行したところ、30 分以内に結果を得ることができなかった。そのため、単純なバックトラッキングアルゴリズムを用いた実験は中止し、代わりにグリーディ彩色アルゴリズムを実装した。

\subsubsection{グリーディ彩色アルゴリズムの実装}
前章\ref{tab:backtracking-only-results}で述べたように, バックトラッキングアルゴリズムを用いて最大クリークを探索した場合, 70ノードを超えるグラフでは結果を得ることができなかった. この問題を解決するために, バックトラッキングアルゴリズムが検証する部分集合の数を削減する手法として, グリーディ彩色アルゴリズムを導入した. グリーディ彩色アルゴリズムとバックトラッキングアルゴリズムを組み合わせた結果は表\ref{tab:backtracking-coloring-results}に示す.

\begin{table}[H]
\centering
\caption{異なるグラフサイズにおける最大クリークサイズと計算時間（グリーディ彩色アルゴリズム）}\label{tab:backtracking-coloring-results}
\begin{tabular}{l|c|r}
\hline
\textbf{グラフサイズ} & \textbf{最大クリークサイズ} & \textbf{計算時間 (s)} \\
\hline
30  & 8  & 0.0046 \\
50  & 9  & 0.0035 \\
70  & 11 & 0.0419 \\
100 & 14 & 0.1951 \\
300 & 18 & 4.8325 \\
\hline
\end{tabular}
\end{table}

この実験において, 組み合わせアルゴリズムは 30, 50, 70, 100 ノードのグラフに対していずれも1秒未満で最大クリークを発見した. それぞれの最大クリークサイズは 8, 9, 11, 14 であった. さらに, 100ノードから300ノードへと規模を拡大した場合でも, 300ノードグラフにおいて最大クリークサイズ18を約4.8 秒で発見することができた.

\section{分析と理論}
\subsection{課題１と課題２:DFSとBFSによる全域木の特徴}
\subsubsection{深さの比較}
The spanning tree returned as a result of traversal via DFS and BFS shows significant difference that emphasizes the main difference between how DFS and BFS work. Comparing the depth of the DFS and BFS spanning tree in table \ref{table:dfs-bfs-depth}, we can see that DFS tends to produce deeper spanning trees no matter dense or sparse. This is due to the core concept of how DFS works which goes as far as possible before backtracking to explore other branches; thus taller or deeper spanning trees. 
\begin{table}[H]
  \centering
  \caption{DFSとBFSにおける全域木の深さの比較}\label{table:dfs-bfs-depth}
  \begin{tabular}{lccc}
  \hline
  ファイル名 & DFSの深さ & BFSの深さ & 差 (DFS-BFS) \\
  \hline
  search\_100s.dat  & 23 & 5 & 18 \\
  search\_100d.dat  & 5  & 3 & 2 \\
  search\_500s.dat  & 37 & 4 & 33 \\
  search\_500d.dat  & 6  & 3 & 3 \\
  search\_1000s.dat & 48 & 3 & 45 \\
  search\_1000d.dat & 9  & 3 & 6 \\
  \hline
  \end{tabular}
\end{table}

On the other hand, executing DFS on sparse graphs tend to produce much deeper spanning trees compared to dense graphs. Sparse graphs contain lower edge density which result in nodes not being well connected to each other. The lower edge count means the algorithms has very few choices of alternative routes to traverse between 1 node and another. This means that there are no shortcuts to get from 1 node to another. The algorithm needs to take a longer route thus creating much deeper spanning trees. Opposite to sparse graphs, dense graphs are much more well connected with bigger edge density. The more interconnected each nodes are to each other creates alternative paths for the algorithm to use to go from a node to another. With a lot of alternatives to choose, the algorithm is able to find shortcuts to shorten the length of path that needs to be taken.
\subsubsection{葉ノード数}
Comparison of the leafs of both spanning trees can be seen in table \ref{table:dfs-bfs-leaf-percent}. While DFS tends to create more leafs than BFS, the difference is extremently minimal and falls within a margin of error of 5\%. The results are negligible and we are able to conclude that neither DFS and BFS have any significant effects on leaf counts.
\begin{table}[H]
  \centering
  \caption{DFSとBFSにおける全域木の葉ノード数の比較（差\%）}\label{table:dfs-bfs-leaf-percent}
  \begin{tabular}{lccc}
  \hline
  ファイル名 & DFSの葉ノード数 & BFSの葉ノード数 & 差 (\%) \\
  \hline
  search\_100s.dat  & 73  & 70  & 4.29\%  \\
  search\_100d.dat  & 96  & 94  & 2.13\%  \\
  search\_500s.dat  & 457 & 453 & 0.88\%  \\
  search\_500d.dat  & 495 & 493 & 0.41\%  \\
  search\_1000s.dat & 947 & 946 & 0.11\%  \\
  search\_1000d.dat & 992 & 993 & -0.10\% \\
  \hline
  \end{tabular}
\end{table}
\subsubsection{最大分岐数}
As explained before, the maximum branching factor calculates the furthest distance a node needs to travel to visit another node. While depth measures from root to the furthest node, the maximum branching factor calculates distance for any 2 nodes. The comparison of spanning trees can be seen in table \ref{table:dfs-bfs-maxbranch}.
\begin{table}[H]
  \centering
  \caption{DFSとBFSにおける全域木の最大分岐数の比較（差\%）}\label{table:dfs-bfs-maxbranch}
  \begin{tabular}{lccc}
  \hline
  ファイル名 & DFS最大分岐数 & BFS最大分岐数 & 差 (\%) \\
  \hline
  search\_100s.dat  & 11   & 14   & -21.43\% \\
  search\_100d.dat  & 66   & 66   & 0.00\%  \\
  search\_500s.dat  & 46   & 49   & -6.12\%  \\
  search\_500d.dat  & 344  & 344  & 0.00\%  \\
  search\_1000s.dat & 107  & 107  & 0.00\%  \\
  search\_1000d.dat & 606  & 606  & 0.00\%  \\
  \hline
  \end{tabular}
\end{table}



\section{まとめ}
\section{発表についての感想}
\section{付録}
本実験に利用したソースコードなどはGitHubにホストされている.
\url{https://github.com/SorataBaka/kosen-5/tree/main/graph-network}
\bibliographystyle{plain}
\bibliography{references}
\end{document}